\section{Mathematical Sentences: Statements, Predicates, and Claims}
\chapterline

Before constructing proofs, we need to inspect the grammar of mathematical language. This is not grammar in the sense of memorizing punctuation rules. It is the physics of how symbols fit together: what kinds of objects exist, which expressions refer to things, which expressions make claims, and which claims can be evaluated as true or false.

\begin{historybox}
For more than two thousand years, formal logic was synonymous with Aristotle's syllogisms (c.~350 BCE in the \emph{Prior Analytics}), which classified arguments into verbal templates like ``All humans are mortal; Socrates is human; therefore Socrates is mortal.'' Scholars had only natural language and rhetorical categories to guide them. But as mathematics advanced toward modern calculus and abstract geometry, Aristotle's grammar broke down completely. It could not express even simple relational deductions, such as Augustus De Morgan's challenge: ``All horses are animals; therefore every head of a horse is a head of an animal.''

Natural language set dangerous syntactic traps. In everyday speech, the little word ``is'' does three entirely different jobs without warning: it asserts identity (``a triangle \emph{is} a three-sided polygon''), declares a property or membership (``seven \emph{is} prime''), or claims universal containment (``a dog \emph{is} an animal''). Conflating these three roles caused centuries of philosophical confusion and made rigorous proofs of infinite processes impossible.

In 1879, the German mathematician Gottlob Frege published his revolutionary \emph{Begriffsschrift} (``Concept Script''). Frege replaced everyday subject-predicate grammar with the mathematical concept of a function. An open condition like ``\ldots is prime'' is a predicate function $P(x)$ that accepts an object and outputs an unambiguous truth value: True or False. This breakthrough gave mathematics its modern taxonomy of terms, predicates, and statements. Yet this taxonomy is purely static; it tells us what a sentence is, but not how variables change or interact across a proof. In Chapter~3 (\emph{Variables Have Roles: Fixed, Arbitrary, Chosen, Bound}), we add the dynamic engine of variable scoping.
\end{historybox}

In mathematics, sentences are divided by this sharp taxonomic line.

\begin{definition}
A \emph{statement} (or \emph{proposition}) is a sentence that is unambiguously either true or false, but not both.
\end{definition}

Consider these examples:
\begin{itemize}
\item ``$7$ is a prime number'' is a true statement.
\item ``$9$ is a prime number'' is a false statement.
\item ``There are infinitely many prime numbers'' is a true statement.
\item ``Every continuous function on $\R$ is differentiable'' is a false statement (the absolute value function $f(x) = |x|$ fails to have a derivative at $x = 0$).
\end{itemize}

A sentence does not need to be easy to evaluate to qualify as a statement. For more than three centuries, Fermat's Last Theorem---the claim that $a^n + b^n = c^n$ has no positive integer solutions for integers $n > 2$---remained unproved. Its truth was unknown to the entire human race until Andrew Wiles and Richard Taylor completed a proof in 1995. Throughout those three centuries, the sentence was not a vague riddle; it was a bona fide statement waiting for a proof.

\subsection{The taxonomy of mathematical expressions}

A common obstacle when transitioning from computational mathematics to proof-based mathematics is conflating nouns, open conditions, assertions, and procedural commands.

\begin{center}
\begin{tikzpicture}[
    font=\small,
    category/.style={draw=#1!70!black, fill=#1!6, rounded corners=3pt, thick, text width=2.7cm, inner sep=6pt, minimum height=4.1cm, align=left},
    header/.style={font=\footnotesize\bfseries\scshape, text=#1!80!black, align=center},
    sample/.style={font=\scriptsize, text=gray!80!black}
]

% Category 1: Terms
\node[category=sourceblue] (col1) at (0, 0) {
    {\centering\textbf{\color{sourceblue!85!black}Mathematical Term}\\\scriptsize(A Noun / Value)\par}
    \vspace{0.2em}
    \rule{2.6cm}{0.4pt}\\[0.35em]
    \scriptsize Evaluates to an object (a number, matrix, or set). Has \emph{no} truth value.\\[0.55em]
    \textbf{Examples:}\\[0.2em]
    $\bullet$~$x^2 + 1$\\
    $\bullet$~$\sin(\theta)$\\
    $\bullet$~$\{1, 4, 9\}$
};

% Category 2: Predicates
\node[category=quotientorange] (col2) at (3.35, 0) {
    {\centering\textbf{\color{quotientorange!85!black}Predicate}\\\scriptsize(Open Template)\par}
    \vspace{0.2em}
    \rule{2.6cm}{0.4pt}\\[0.35em]
    \scriptsize Contains free variable slots. Truth value is suspended until slots are filled.\\[0.55em]
    \textbf{Examples:}\\[0.2em]
    $\bullet$~$x^2 - 1 = 0$\\
    $\bullet$~$n$ is prime\\
    $\bullet$~$A \subseteq B$
};

% Category 3: Statements
\node[category=targetgreen] (col3) at (8.15, 0) {
    {\centering\textbf{\color{targetgreen!85!black}Statement}\\\scriptsize(Closed Claim)\par}
    \vspace{0.2em}
    \rule{2.6cm}{0.4pt}\\[0.35em]
    \scriptsize Completely closed assertion. Definitively either True or False.\\[0.55em]
    \textbf{Examples:}\\[0.2em]
    $\bullet$~$7$ is prime (\textbf{T})\\
    $\bullet$~$2 + 2 = 5$ (\textbf{F})\\
    $\bullet$~$\forall x \in \R,\, x^2 \ge 0$ (\textbf{T})
};

% Category 4: Type Errors
\node[category=trapred] (col4) at (11.5, 0) {
    {\centering\textbf{\color{trapred!85!black}Type Error}\\\scriptsize(Ill-Formed)\par}
    \vspace{0.2em}
    \rule{2.6cm}{0.4pt}\\[0.35em]
    \scriptsize Syntactically incompatible types. Meaningless to evaluate.\\[0.55em]
    \textbf{Examples:}\\[0.2em]
    $\bullet$~$5 \subseteq \R$\\
    $\bullet$~$x + 2 = \text{circle}$\\
    $\bullet$~$\sin \land \cos$
};

% Transformation arrows from Predicate to Statement in dedicated corridor
\draw[-{Latex[length=2.2mm]}, thick, quotientorange!80!black] ([yshift=0.7cm]col2.east) -- node[above=2pt, font=\tiny\bfseries, text=quotientorange!90!black]{Instantiate ($x \leftarrow 1$)} ([yshift=0.7cm]col3.west);
\draw[-{Latex[length=2.2mm]}, thick, quotientorange!80!black] ([yshift=-0.7cm]col2.east) -- node[below=2pt, font=\tiny\bfseries, text=quotientorange!90!black]{Quantify ($\exists x$)} ([yshift=-0.7cm]col3.west);

\end{tikzpicture}
\end{center}

\subsection{Predicates: sentences with empty sockets}

An expression such as $x^2 - 1 = 0$ looks like a statement, but it is not. You cannot say whether it is true or false until you know what $x$ is. If $x = 1$, the sentence becomes $1^2 - 1 = 0$, which is true. If $x = 3$, the sentence becomes $3^2 - 1 = 0$, which is false.

\begin{definition}
A \emph{predicate} is a sentence containing one or more variables that becomes a statement whenever each variable is replaced by a specific element from a specified domain, or bound by a quantifier.
\end{definition}

In computer science, a predicate is a function that returns a boolean value:
\begin{center}
\texttt{isEven : Integer $\to$ \{True, False\}}
\end{center}
Until you pass an input into \texttt{isEven(n)}, the function call produces no boolean output. In mathematics, we write $P(x)$ to indicate that $P$ is a predicate waiting for an input $x$.

\subsection{The domain changes everything}

When you work with a predicate, the mathematical universe from which variables may be chosen---the \emph{domain of discourse}---is not an optional footnote. It governs the truth of the sentence.

Recall our standard number systems:
\begin{itemize}
\item $\N = \{1, 2, 3, \ldots\}$: the natural numbers (positive integers).
\item $\Z = \{\ldots, -2, -1, 0, 1, 2, \ldots\}$: the integers.
\item $\Q = \{p/q : p, q \in \Z, q \neq 0\}$: the rational numbers (fractions).
\item $\R$: the real numbers (including rationals and irrationals like $\sqrt{2}$ and $\pi$).
\item $\C = \{a + bi : a, b \in \R, i^2 = -1\}$: the complex numbers.
\end{itemize}

Now examine the predicate $P(x) : x^2 < 2$. Its truth depends entirely on where $x$ is allowed to live:
\begin{itemize}
\item Over the natural numbers $\N$, the statement ``$\exists x \in \N, x^2 < 2$'' has only a single witness: $x = 1$.
\item Over the integers $\Z$, the solutions are $x \in \{-1, 0, 1\}$.
\item Over the rational numbers $\Q$, there are infinitely many solutions, yet the boundary points $-\sqrt{2}$ and $\sqrt{2}$ do not exist inside $\Q$.
\item Over the real numbers $\R$, the solution set is the open interval $(-\sqrt{2}, \sqrt{2})$.
\item Over the complex numbers $\C$, the inequality $x^2 < 2$ is not even standardly defined, because the complex numbers do not possess an ordering compatible with arithmetic.
\end{itemize}

\begin{trapbox}
Never leave a variable floating without a declared domain. Writing $x^2 \ge 0$ without specifying the domain is an invitation to confusion. Over $\R$, it is true; over $\C$, choosing $x = i$ gives $i^2 = -1 < 0$.
\end{trapbox}

\subsection{Claims, theorems, and mathematical etiquette}

When mathematicians state assertions they intend to establish, they use conventional labels to communicate the structural role of the result:

\begin{itemize}
\item \textbf{Theorem}: A central, landmark result that anchors a theory (e.g., the Pythagorean Theorem, the Fundamental Theorem of Algebra).
\item \textbf{Proposition}: A solid, standalone result of intermediate significance, worth recording in its own right.
\item \textbf{Lemma}: A stepping-stone result. Lemmas often do the heavy algebraic lifting so that the main theorem can be stated cleanly.
\item \textbf{Corollary}: A direct and immediate consequence of an established theorem, requiring minimal additional effort to verify.
\item \textbf{Conjecture}: A statement believed to be true based on strong evidence or heuristic arguments, but which currently lacks an inescapable proof.
\end{itemize}

These terms are stylistic rather than formal. One author's lemma might be celebrated by later generations as a profound theorem (such as Zorn's Lemma or Gauss's Lemma).

\subsection{What formal grammar unlocked---and its operational limits}

Frege's functional taxonomy unlocked an unprecedented clarity: it permanently separated the object of study (nouns/terms) from the assertions made about it (statements), while providing a rigorous home for open equations through predicates. It protects mathematicians from catastrophic category mistakes, such as trying to assert that a number is a subset, or testing whether an open formula is true without binding its variables.

Yet this grammatical taxonomy has an honest operational limitation: \textbf{it is completely static}. It tells you what an expression is, but it says nothing about how variables behave over the course of a multi-paragraph proof. If a sentence contains an $x$, who chooses its value? Is $x$ an immovable constant given by the problem, an adversarial input you have no power to alter, an object you are free to construct, or a private dummy counter that will vanish at the end of the line? 

Without clear rules governing variable lifetime, visibility, and permissions, even grammatically perfect sentences can lead to circular or invalid proofs. In Chapter~3 (\emph{Variables Have Roles: Fixed, Arbitrary, Chosen, Bound}), we move from static grammar to dynamic execution, mastering the four fundamental roles of variables and the iron laws of lexical scope.

\subsection{Exercises}

Classify each expression below as a term, a predicate, a statement, an instruction, or a grammatical type error:
\begin{enumerate}
\item $2x + 5 = 11$.
\item $2(3) + 5 = 11$.
\item $\cos^2(\theta) + \sin^2(\theta)$.
\item For every real number $x$, $x^2 + 1 > 0$.
\item Let $G$ be a finite group.
\item $3 \in \{1, 2, 3, 4\}$.
\item $\{3\} \in \{1, 2, 3, 4\}$.
\item $\{3\} \subseteq \{1, 2, 3, 4\}$.
\item $\sqrt{-1} \le 4$.
\end{enumerate}

For each of the following predicates, describe two different domains where the predicate yields different sets of solutions:
\begin{enumerate}[resume]
\item $Q(x) : 2x = 7$.
\item $R(x) : x^2 = 2$.
\item $S(x) : x \ge 0$.
\end{enumerate}

\newpage
